% problem_id: S014-lurie-ultracategories-theorem-2-4-2-makkai
% source_id: S014 (Jacob Lurie, Ultracategories)
% source_file: lurie-conceptual.pdf
% statement_locator: printed p. 27
% proof_locator: printed p. [29]
% representation: PDF; transcribed from rendered page images (never from the text layer)
% status: complete (statement + author proof verbatim + reason + locators)
%
%% problem_id: L3
%% source_locator: J. Lurie, "Ultracategories", Theorem 2.4.2 (Makkai)
%%   PDF: /disks/sata1/yupeng/human-proof-corpus/source-cache/registry-sources/lurie/lurie-conceptual.pdf
%%   statement page: printed p. 27 (PDF page 27 — printed and PDF page numbers coincide in this file)
%%   proof page:     printed p. 29 (PDF page 29 — the whole proof is contained on this one page)
%%   transcription is from 150 dpi / 300 dpi pdftoppm renderings of the original pages that were read as images.
%%   The PDF text layer was used ONLY for page location, never as the source of wording
%%   (it corrupts math symbols: e.g. the ACL "∈" of the text layer is unreliable, "α" is printed as a broken glyph).
%% Notes on printing oddities are given inline as %% NOTE; nothing has been silently corrected.

\documentclass[11pt]{article}
\usepackage{amsmath,amssymb,amsthm}
\usepackage[utf8]{inputenc}
\usepackage[T1]{fontenc}

\newcommand{\E}{\mathcal{E}}          % the regular / exact category, script E
\newcommand{\C}{\mathcal{C}}          % the pretopos, script C
\newcommand{\Fun}{\operatorname{Fun}}
\newcommand{\Funreg}{\Fun^{\mathrm{reg}}}
\newcommand{\Mod}{\operatorname{Mod}}
\newcommand{\Set}{\mathbf{Set}}
\newcommand{\ev}{\mathrm{ev}}
\newcommand{\Shv}{\operatorname{Shv}}

\begin{document}

%% =====================================================================
%% --- page 27 ---  (statement page)
%% =====================================================================

\noindent\textbf{Theorem 2.4.2 (Makkai).} \emph{Let $\E$ be a small exact category (Definition A.2.6). Then the functor $\ev : \E \to \Fun(\Funreg(\E,\Set),\Set)$ is a fully faithful embedding, whose essential image consists of those functors $F : \Funreg(\E,\Set) \to \Set$ which preserve small products and small filtered colimits.}

%% --- end of page 27 statement ---

%% ---------------------------------------------------------------------
%% CONTEXT (not part of Theorem 2.4.2; transcribed because the proof of
%% Theorem 2.4.2 refers to it twice). Also printed on page 27.
%% ---------------------------------------------------------------------

\medskip
\noindent\textbf{Proposition 2.4.3.} \emph{Let $\E$ be a small regular category. Then there exists a small pretopos $\C$ and a fully faithful regular functor $h : \E \to \C$ such that precomposition with $h$ induces an equivalence of categories $\Mod(\C) \subseteq \Funreg(\C,\Set) \to \Funreg(\E,\Set)$. Moreover, if the category $\E$ is exact, then an object $C \in \C$ belongs to the essential image of $f$ if and only if the evaluation functor}
\[
\ev_C : \Mod(\C) \to \Set \qquad M \mapsto M(C)
\]
\emph{commutes with finite products.}

%% NOTE (printed oddity, left as printed): in the second sentence of Proposition 2.4.3
%% the essential image is said to be "of $f$", although the functor introduced in the
%% proposition is called $h$. Presumably a typo for $h$; not corrected here.
%% NOTE (printed oddity, left as printed): the displayed chain "$\Mod(\C) \subseteq \Funreg(\C,\Set) \to \Funreg(\E,\Set)$"
%% is printed with the symbol $\subseteq$ between the first two terms. Left as printed.

%% =====================================================================
%% --- page 29 ---  (proof page; the entire proof of Theorem 2.4.2 sits on this page)
%% =====================================================================

\medskip
\noindent\emph{Proof of Theorem 2.4.2.} Let $\E$ be a small exact category. According to Theorem 2.4.1, the construction $E \mapsto \ev_E$ induces a fully faithful embedding $\ev : \E \hookrightarrow \Fun(\Funreg(\E,\Set),\Set)$. It will therefore suffice to show that if $F : \Funreg(\E,\Set) \to \Set$ is a functor which preserves small products and small filtered colimits, then there is a natural isomorphism $F \simeq \ev_E$ for some $E \in \E$ (the converse is clear, since small products and small filtered colimits in the category $\Funreg(\E,\Set)$ are computed pointwise). Note that we can use Proposition 1.4.9 to endow $F$ with the structure of an ultrafunctor (where $\Funreg(\E,\Set)$ and $\Set$ are equipped with the categorical ultrastructures of Example 1.3.8).

Choose a small pretopos $\C$ and a regular functor $h : \E \to \C$ satisfying the requirements of Proposition 2.4.3. Then precomposition with $h$ induces an equivalence of ultracategories $H : \Mod(\C) \to \Funreg(\E,\Set)$ (Remark 2.4.5). Consequently, the composite functor
\[
\Mod(\C) \xrightarrow{\ H\ } \Funreg(\E,\Set) \xrightarrow{\ F\ } \Set
\]
admits an ultrastructure, and is therefore given by evaluation on some object $C \in \C$ (Theorem 2.3.1). It will therefore suffice to show that $C$ belongs to the essential image of $h$. This follows from Proposition 2.4.3, since the functor $F \circ H$ commutes with finite products. \hfill$\square$

%% --- end of page 29 proof ---

%% =====================================================================
%% TRANSCRIPTION SCOPE
%% =====================================================================
%% The proof of Theorem 2.4.2 does NOT occupy page 28. Page 28 (printed p. 28)
%% contains the proof of Proposition 2.4.3 (its clauses ($\ast$) and ($\ast'$)), which is a
%% separate result and is not part of the proof of Theorem 2.4.2; it was rendered and
%% read for context but is deliberately not transcribed here.
%% No portion of the statement or of the proof was illegible; no [ILLEGIBLE: ...]
%% marker was required.

\end{document}
